\section{Introduction}

Weekly influenza hospitalization forecasting sits at an awkward intersection
between mechanistic epidemiology and operational prediction. On one hand,
compartmental epidemic models remain attractive because they offer interpretable
latent states, controllable inductive bias, and explicit constraints such as
non-negativity and approximate mass conservation. On the other hand, practical
forecast targets such as hospitalization rate are not observed directly as
canonical infectious prevalence. A model can therefore fail not only because
its latent disease dynamics are misspecified, but also because its
\emph{observation semantics} are mismatched to the public-health quantity being
predicted.

This mismatch is especially relevant for a proof-of-concept benchmark based on
a single FluSurv-NET season. In such a setting, one should not expect a single
model family to dominate across all age strata. Different age groups may differ
in latency, effective susceptibility recycling, reporting delay, and the extent
to which observed hospitalization rates align with latent infectious burden.
Accordingly, the central scientific question is not simply which model wins on
one leaderboard, but rather:
\begin{quote}
Can a constrained, fully programmatic structure-discovery loop identify when
simple mechanistic baselines are sufficient and when alternative latent or
observation structures become useful?
\end{quote}

This repository deliberately answers that question without placing a language
model inside the fitting loop. The current benchmark treats the ``agentic''
component as a propose-fit-verify-refine search procedure over a small grammar
of epidemic structures and observation maps. Candidate models are generated
programmatically, screened for biological and structural validity, fitted with
multiple random restarts, scored on validation and rolling-origin behavior, and
refined through beam search. This design makes the non-LLM control condition
strong enough to support a later LLM-based extension without confusing search
automation with semantic reasoning.

The benchmark evolved in three substantive steps. First, it established a
reproducible non-LLM forecasting core spanning deterministic, probabilistic,
and fractional SEIR baselines together with constrained structure discovery.
Second, it strengthened the comparison set with fair hospitalization-aware
manual baselines and benchmark-level conformal uncertainty post-processing.
Third, it elevated observation structure to a first-class discovery object by
allowing delayed infectious observation and hospitalization-aware observation
maps inside the search grammar. This last step matters because the target is a
hospitalization-rate series: if discovery is allowed to vary only latent
compartments while the observation map remains effectively fixed, then the
comparison between hand-designed and discovered models is incomplete.

The main empirical outcome is now clear. The benchmark does \emph{not} support
a single globally dominant model family. Instead, it supports
\emph{age-aware model selection under structural and observation uncertainty}.
In the current five-seed observation-aware benchmark:
\begin{itemize}[leftmargin=1.5em]
\item constrained discovery is a stable success case for the \texttt{0--4 yr}
series;
\item stronger manual baselines dominate or match discovery for
\texttt{Overall}, \texttt{18--49 yr}, and \texttt{50--64 yr};
\item delayed infectious observation emerges repeatedly in children and older
adults;
\item direct hospitalization observation is not selected by discovery in the
current single-season benchmark; and
\item removing the age prior does not change the selected structures,
observation maps, delays, or aggregate discovery errors.
\end{itemize}

These results turn the repository into a strong non-LLM experimental core for a
subsequent LLM phase. Because the current benchmark already contains fair manual
baselines, observation-aware discovery, multi-seed robustness analysis,
benchmark-level uncertainty calibration, and no-age-prior ablations, the next
step can focus narrowly on whether an LLM improves candidate proposal and
critique efficiency rather than on whether the underlying fitting and
evaluation pipeline is sound.
